\clearpage
\section{Metodyka badań i opracowanie algorytmu sterowania}
\label{sec:metodyka}

\subsection{Koncepcja eksperymentu}
\label{subsec:koncepcja}

Badanie zaprojektowano jako kontrolowany eksperyment porównawczy, w którym agent uczony metodą uczenia ze wzmocnieniem oraz dwie klasyczne metody referencyjne sterują tą samą siecią drogową i są oceniane na tych samych scenariuszach ruchu. Nadrzędnym założeniem metodycznym było oddzielenie danych wykorzystywanych do uczenia i strojenia modelu od danych służących do końcowej oceny.

Procedurę badawczą podzielono na następujące etapy:
\begin{enumerate}
    \item przygotowanie modelu skrzyżowania w środowisku symulacyjnym SUMO;
    \item wygenerowanie scenariuszy ruchu odpowiadających różnym poziomom natężenia;
    \item podział danych na zbiór treningowy, walidacyjny i testowy;
    \item dobór parametrów czasowych metod referencyjnych na podstawie scenariuszy treningowych;
    \item trening modeli DQN na przygotowanych danych treningowych;
    \item wyznaczenie odpowiedniej liczby kroków treningowych oraz porównanie dwóch sposobów generowania danych treningowych;
    \item dobór hiperparametrów metodą zmiany pojedynczego parametru;
    \item zamrożenie wybranej konfiguracji jako modelu końcowego;
    \item końcowa ocena modelu na odrębnym zbiorze testowym;
    \item porównanie wyników agenta DQN z wynikami metod klasycznych.
\end{enumerate}

Wybór i strojenie modelu DQN oraz dobór parametrów metod referencyjnych realizowano wyłącznie na danych treningowych i walidacyjnych. Zbiór testowy służył do porównania zamrożonych konfiguracji metod i nie stanowił formalnego kryterium wyboru DQN. Chronologię jego wykorzystania oraz wynikające z niej ograniczenie omówiono w podrozdziale~\ref{subsec:ograniczenia-wiarygodnosci}. Schemat całego procesu badawczego przedstawiono na rysunku~\ref{fig:proces-badawczy}.

\begin{figure}[!htbp]
    \centering
    \includegraphics[width=0.65\textwidth]{proces_badawczy.png}
    \caption{Schemat procedury badawczej z rozdzieleniem danych treningowych, walidacyjnych i testowych.}
    \label{fig:proces-badawczy}
\end{figure}

\subsection{Charakterystyka analizowanego skrzyżowania}
\label{subsec:skrzyzowanie}

Do badania wybrano skrzyżowanie ulicy Belgradzkiej z aleją Komisji Edukacji Narodowej, położone w dzielnicy Ursynów w Warszawie. Skrzyżowanie wybrano jako stosunkowo prosty pod względem geometrii, lecz istotnie obciążony układ drogowy, w którym aleja Komisji Edukacji Narodowej stanowi kierunek główny. Położenie skrzyżowania w skali dzielnicy przedstawiono na rysunku~\ref{fig:mapa-ursynow}, a jego bezpośrednie otoczenie na rysunkach~\ref{fig:mapa-osm} oraz~\ref{fig:satelita}.

\begin{figure}[!htbp] \centering
    \includegraphics[width=0.75\linewidth]{map-2.png}
    \caption{Położenie analizowanego skrzyżowania na tle warszawskiego Ursynowa. Skrzyżowanie zaznaczono czerwonym okręgiem. Źródło: OpenStreetMap.}
    \label{fig:mapa-ursynow}
\end{figure}

\begin{figure}[!htbp] \centering
    \includegraphics[width=0.75\linewidth]{map.png}
    \caption{Analizowane skrzyżowanie na mapie OpenStreetMap.}
    \label{fig:mapa-osm}
\end{figure}

\begin{figure}[!htbp] \centering
    \includegraphics[width=0.75\linewidth]{satelita.png}
    \caption{Widok satelitarny analizowanego skrzyżowania. Źródło: Google Maps.}
    \label{fig:satelita}
\end{figure}

Skrzyżowanie ma cztery wloty, oznaczane w dalszej części pracy jako północny, południowy, wschodni oraz zachodni. Nazwy te odnoszą się do położenia wlotu względem środka skrzyżowania, a nie do kierunku jazdy. Model dopuszcza ruch z każdego wlotu we wszystkich dostępnych kierunkach, z wyłączeniem zawracania. Liczba pasów na poszczególnych wlotach jest zróżnicowana: wlot południowy, z którego pojazdy poruszają się w kierunku północnym, obejmuje cztery pasy, natomiast pozostałe trzy wloty po trzy pasy, co daje łącznie trzynaście pasów wlotowych objętych sterowaniem.

Program sygnalizacji obejmuje trzy fazy z sygnałem zielonym, rozdzielone fazami żółtymi. Pierwsza faza obsługuje jazdę na wprost oraz skręt w prawo na alei Komisji Edukacji Narodowej, jednocześnie z obu kierunków tej osi. Druga faza realizuje skręt w lewo z alei Komisji Edukacji Narodowej, chroniony osobnym sygnałem i rozdzielony w czasie od kolizyjnego strumienia jadącego na wprost z kierunku przeciwnego. Trzecia faza obsługuje ulicę Belgradzką, obejmując jazdę na wprost, skręt w prawo oraz skręt w lewo realizowany warunkowo, to jest przy ustąpieniu pierwszeństwa pojazdom nadjeżdżającym z przeciwka. Różnica w sposobie obsługi lewoskrętów wpływa na przepustowość tych relacji i dlatego jest istotna dla interpretacji wyników.

W modelu odwzorowano geometrię skrzyżowania, liczbę i przeznaczenie pasów ruchu, relacje kolizyjne oraz program sygnalizacji świetlnej. Model zbudowano na podstawie danych OpenStreetMap, a następnie skorygowano na potrzeby eksperymentu. Odwzorowane skrzyżowanie w środowisku SUMO przedstawiono na rysunku~\ref{fig:symulacja}.

\begin{figure}[!htbp] \centering
    \includegraphics[width=0.8\linewidth]{symulacja.png}
    \caption{Model analizowanego skrzyżowania w środowisku symulacyjnym SUMO.}
    \label{fig:symulacja}
\end{figure}

Zgodnie z ograniczeniami zakresu przedstawionymi w podrozdziale~\ref{subsec:zalozenia-ograniczenia}, model nie obejmuje faz ani ograniczeń czasowych związanych z obsługą przejść dla pieszych.

\subsection{Sformułowanie problemu jako zadania uczenia ze wzmocnieniem}
\label{subsec:sformulowanie-rl}

\subsubsection{Agent i środowisko}
\label{subsubsec:agent-srodowisko}

Problem sterowania sygnalizacją sformułowano jako zadanie uczenia ze wzmocnieniem. Agent odpowiada za wybór fazy sygnalizacji świetlnej, natomiast środowisko stanowi symulacja ruchu drogowego realizowana w SUMO. Komunikację między agentem a symulatorem zapewnia interfejs TraCI, który umożliwia wykonywanie kolejnych kroków symulacji, odczyt stanu ruchu oraz zmianę sygnalizacji.

Warstwę pośredniczącą między symulatorem a algorytmem uczenia stanowi biblioteka \texttt{sumo-rl}, udostępniająca symulację jako środowisko zgodne z interfejsem Gymnasium \cite{alegre2019sumorl}. Na jej podstawie przygotowano własną warstwę integracyjną. Jej architekturę i implementację opisano w rozdziałach~\ref{sec:architektura} oraz~\ref{sec:implementacja}.

\subsubsection{Przestrzeń obserwacji}
\label{subsubsec:obserwacja}

Przestrzeń obserwacji jest ciągła i ma postać \texttt{Box(0.0, 1.0, (30,), float32)}, czyli wektora trzydziestu wartości rzeczywistych z przedziału od 0 do 1. Wykorzystano domyślną reprezentację stanu udostępnianą przez bibliotekę \texttt{sumo-rl} \cite{alegre2019sumorl}, na którą składają się cztery grupy cech:
\begin{itemize}
    \item \emph{kodowanie aktywnej fazy} (3 wartości): wektor zero-jedynkowy typu one-hot wskazujący, która z trzech faz zielonych jest aktualnie realizowana;
    \item \emph{znacznik minimalnego czasu zielonego} (1 wartość): zmienna binarna przyjmująca wartość 1, jeżeli od ostatniej zmiany fazy upłynął już minimalny czas zielony powiększony o czas trwania fazy żółtej, a w przeciwnym razie wartość 0;
    \item \emph{gęstość pasów} (13 wartości): dla każdego z trzynastu pasów wlotowych liczba pojazdów znajdujących się na pasie podzielona przez maksymalną liczbę pojazdów, które mogłyby się na nim zmieścić;
    \item \emph{zapełnienie kolejki} (13 wartości): dla każdego pasa wlotowego liczba pojazdów zatrzymanych, czyli poruszających się z prędkością poniżej 0,1~m/s, podzielona przez tę samą pojemność pasa.
\end{itemize}

Pojemność pasa wyznaczana jest jako jego długość podzielona przez sumę minimalnego odstępu między pojazdami oraz średniej długości pojazdu na tym pasie. Obie ostatnie grupy cech są ograniczone od góry wartością 1, co zapewnia normalizację całego wektora obserwacji do wspólnego zakresu i ułatwia uczenie sieci neuronowej.

Przyjęta reprezentacja opisuje zatem jednocześnie stan sygnalizacji, dopuszczalność zmiany fazy oraz rozkład obciążenia ruchem na poszczególnych pasach wlotowych. Rozdzielenie gęstości i zapełnienia kolejki pozwala odróżnić sytuację, w której pojazdy znajdują się na pasie i płynnie przejeżdżają, od sytuacji, w której pojazdy stoją w kolejce.

\subsubsection{Przestrzeń akcji}
\label{subsubsec:akcje}

Przestrzeń akcji jest dyskretna i ma postać \texttt{Discrete(3)}. Każda akcja odpowiada wskazaniu fazy zielonej, która ma obowiązywać w kolejnym oknie decyzyjnym \cite{alegre2019sumorl}. Trzy dostępne akcje odpowiadają trzem fazom zielonym zdefiniowanym w programie sygnalizacji, opisanym w podrozdziale~\ref{subsec:skrzyzowanie}.

Agent nie steruje bezpośrednio sygnałami na poszczególnych pasach ani nie ustala długości fazy w sposób dowolny. Wskazuje jedynie fazę docelową, natomiast egzekwowanie ograniczeń bezpieczeństwa pozostaje po stronie środowiska. Obowiązują przy tym następujące reguły:
\begin{itemize}
    \item jeżeli agent wskaże fazę identyczną z aktualnie realizowaną, sygnalizacja kontynuuje bieżącą fazę bez żadnej zmiany;
    \item jeżeli agent wskaże fazę inną niż aktualna, lecz od ostatniej zmiany nie upłynął jeszcze minimalny czas zielony powiększony o czas fazy żółtej, żądanie zmiany jest ignorowane i bieżąca faza jest kontynuowana;
    \item jeżeli agent wskaże fazę inną niż aktualna, a warunek czasowy jest spełniony, środowisko wprowadza fazę żółtą jako przejście, a dopiero po jej zakończeniu włącza wskazaną fazę zieloną.
\end{itemize}

Faza żółta jest generowana automatycznie na podstawie porównania stanu sygnałów fazy bieżącej i docelowej, co zapewnia wygaszanie tylko tych relacji, które tracą prawo przejazdu. Mechanizm ten sprawia, że agent nie może przełączać sygnałów natychmiastowo ani z pominięciem czasu międzyzielonego. W eksperymencie przyjęto czas trwania fazy żółtej równy 3~s oraz minimalny czas zielony równy 5~s. Te same wartości obowiązywały podczas treningu, walidacji oraz końcowej oceny na zbiorze testowym, dzięki czemu warunki działania agenta były identyczne na wszystkich etapach eksperymentu.

\subsubsection{Funkcja nagrody}
\label{subsubsec:nagroda}

Wykorzystano funkcję nagrody opartą na zmianie skumulowanego czasu oczekiwania pojazdów \cite{alegre2019sumorl}. W konfiguracji nosi ona nazwę \texttt{diff-waiting-time}. W chwili decyzyjnej $t$ wyznaczana jest wielkość
\begin{equation}
    W_t = \frac{1}{100} \sum_{l \in L} w_{l,t},
    \label{eq:waiting-total}
\end{equation}
gdzie $L$ oznacza zbiór pasów wlotowych sterowanych przez sygnalizację, a $w_{l,t}$ część skumulowanego czasu oczekiwania pojazdów znajdujących się na pasie $l$, przypisaną temu pasowi. Skumulowany czas oczekiwania pojedynczego pojazdu raportowany jest przez SUMO i obejmuje czas, w którym pojazd poruszał się z prędkością poniżej progu zatrzymania. Jeżeli pojazd zmienił pas w obrębie sterowanego obszaru, wykorzystana biblioteka rozdziela tę wielkość między kolejno zajmowane pasy, dzięki czemu ten sam czas oczekiwania nie jest uwzględniany wielokrotnie. Czynnik $\frac{1}{100}$ skaluje sumę do mniejszego zakresu liczbowego. Nagroda przekazywana agentowi jest różnicą między wartością z poprzedniej i bieżącej chwili decyzyjnej:
\begin{equation}
    r_t = W_{t-1} - W_t.
    \label{eq:reward}
\end{equation}

Nagroda jest zatem dodatnia, gdy łączny czas oczekiwania na skrzyżowaniu maleje, a ujemna, gdy rośnie. Agent nie jest uczony bezpośrednio minimalizacji wartości bezwzględnej czasu oczekiwania, lecz podejmowania decyzji prowadzących do jego redukcji względem stanu poprzedniego. Taka konstrukcja daje gęsty sygnał uczący, dostępny po każdej decyzji, co jest korzystne w zadaniu o długim horyzoncie czasowym.

Wybór tej funkcji wynikał z jej bezpośredniego powiązania z jedną z podstawowych metryk oceny jakości sterowania oraz z dostępności sprawdzonej implementacji w wykorzystanej bibliotece. Należy jednak podkreślić, że funkcja nagrody reprezentuje tylko jeden z możliwych sposobów zdefiniowania celu sterowania. Nie uwzględnia ona wprost długości kolejek, przepustowości ani częstotliwości przełączeń sygnalizacji, co należy brać pod uwagę przy interpretacji wyników w pozostałych metrykach.

\subsubsection{Interwał decyzyjny}
\label{subsubsec:interwal}

Agent podejmuje decyzje w stałych odstępach czasu symulacji wynoszących 5~s. Pomiędzy kolejnymi decyzjami symulator wykonuje kroki o długości jednej sekundy, w trakcie których egzekwowane są ograniczenia dotyczące faz oraz rejestrowany jest stan ruchu.

Dobór interwału decyzyjnego stanowi kompromis między częstotliwością reagowania a stabilnością sterowania. Zbyt krótki interwał zwiększa liczbę decyzji przypadających na epizod oraz ryzyko częstych przełączeń sygnalizacji, natomiast zbyt długi ogranicza zdolność agenta do reagowania na zmiany sytuacji ruchowej. Wartość 5~s pozwala agentowi wielokrotnie ocenić sytuację w trakcie trwania jednej fazy, a jednocześnie nie prowadzi do przełączania sygnalizacji częściej, niż dopuszczają to ograniczenia minimalnego czasu zielonego oraz fazy żółtej.

\subsection{Wybór algorytmu}
\label{subsec:wybor-algorytmu}

Punktem wyjścia był klasyczny, tablicowy Q-Learning, opisany w podrozdziale~\ref{subsec:qlearning}. Podejście to odrzucono, ponieważ zastosowanie go do ciągłej, 30-wymiarowej obserwacji wymagałoby dyskretyzacji prowadzącej do niepraktycznie dużej liczby potencjalnych stanów.

Wybrano algorytm Deep Q-Network, który aproksymuje funkcję wartości akcji za pomocą sieci neuronowej \cite{mnih2013playing}. Odpowiada on przyjętemu sformułowaniu zadania, obejmującemu ciągłą obserwację i trzy dyskretne akcje. Wykorzystano implementację ze Stable-Baselines3 oraz politykę \texttt{MlpPolicy}, odpowiednią dla wektorowej reprezentacji obserwacji.

Należy zaznaczyć, że wybór algorytmu DQN nie oznacza, iż jest on rozwiązaniem obiektywnie najlepszym dla rozpatrywanego problemu. Stanowi on natomiast naturalny i uzasadniony wybór przy przyjętym sformułowaniu zadania oraz zakresie pracy.

\subsection{Metody referencyjne}
\label{subsec:metody-referencyjne}

Aby ocena agenta DQN była osadzona w realnym punkcie odniesienia, przygotowano dwie klasyczne metody sterowania. Obie realizują ten sam zbiór faz sygnalizacji co agent, dzięki czemu różnice w wynikach wynikają ze sposobu podejmowania decyzji, a nie z odmiennej organizacji ruchu na skrzyżowaniu.

\subsubsection{Sterowanie stałoczasowe}
\label{subsubsec:ref-staloczasowe}

Sterowanie stałoczasowe stanowi podstawowy punkt odniesienia i w SUMO jest realizowane jako statyczna logika sygnalizacji \cite{sumoTrafficLights}. Czasy faz zielonych wyznaczono proporcjonalnie do natężenia na pasie krytycznym każdej fazy, oszacowanego na podstawie zbioru treningowego. Wynoszą one 14~s, 29~s oraz 29~s, a rozdzielające je fazy żółte trwają po 3~s, co daje cykl o długości 81~s.

\subsubsection{Sterowanie akomodacyjne}
\label{subsubsec:ref-akomodacyjne}

Sterowanie akomodacyjne typu gap-based stanowi bardziej wymagający punkt odniesienia. Może wydłużać aktywną fazę, dopóki detektory rejestrują napływ pojazdów, w granicach minimalnego i maksymalnego czasu zielonego \cite{sumoTrafficLights, NAP22097}. Przyjęto minimum równe 5~s, takie samo jak dla agenta, oraz maksima odpowiadające około półtorakrotności czasów zielonych programu stałoczasowego. Metoda reaguje na ruch według ustalonych reguł, lecz nie uczy się polityki na podstawie doświadczenia.

\subsection{Metryki oceny}
\label{subsec:metryki}

Metody porównywane są według kilku metryk, ponieważ pojedyncza wielkość nie opisuje wszystkich aspektów jakości sterowania. Dane pochodzą z dwóch źródeł: z pliku \texttt{tripinfo.xml} generowanego przez SUMO po zakończeniu symulacji, zawierającego podsumowanie każdej zrealizowanej podróży, oraz z danych rejestrowanych przez interfejs TraCI w każdej sekundzie symulacji. Dla wszystkich metod zastosowano ten sam sposób zbierania i agregacji danych.

\subsubsection{Średnie opóźnienie}
\label{subsubsec:metryka-opoznienie}

Średnie opóźnienie wyznaczane jest jako średnia arytmetyczna wartości \texttt{timeLoss} ze wszystkich zakończonych podróży \cite{sumo_tripinfo_docs}. Wielkość ta określa czas utracony przez pojazd wskutek poruszania się z prędkością niższą od prędkości, z jaką mógłby jechać w warunkach swobodnych. Obejmuje zatem zarówno pełne zatrzymania, jak i zwolnienia wynikające z hamowania oraz przyspieszania.

Metryka wyrażana jest w sekundach na pojazd, a wartości niższe są korzystniejsze. Ze względu na szerokie ujęcie strat czasu metryka ta pełniła wiodącą rolę przy porównywaniu modeli na etapie walidacji.

\subsubsection{Średni czas oczekiwania}
\label{subsubsec:metryka-oczekiwanie}

Wartość \texttt{waitingTime} oznacza czas, w którym pojazd poruszał się poniżej progu uznawanego przez SUMO za zatrzymanie \cite{sumo_tripinfo_docs}. Średni czas oczekiwania jest średnią arytmetyczną tej wartości ze wszystkich zakończonych podróży.

Metryki tej nie należy utożsamiać z opóźnieniem. Opóźnienie obejmuje wszystkie straty czasu, w tym te wynikające ze zwolnienia bez zatrzymania, natomiast czas oczekiwania dotyczy jedynie postoju. Pojazd, który wielokrotnie zwalnia, lecz nigdy się nie zatrzymuje, wykazuje niezerowe opóźnienie przy zerowym czasie oczekiwania. Metryka wyrażana jest w sekundach na pojazd, a wartości niższe są korzystniejsze. Dodatkowo rejestrowano wartość maksymalną oraz kwantyl rzędu 0,95 czasu oczekiwania, opisujące sytuację pojazdów obsłużonych najgorzej.

\subsubsection{Średnia długość kolejki}
\label{subsubsec:metryka-kolejka}

Długość kolejki wyznaczana jest na podstawie danych rejestrowanych w każdej sekundzie symulacji. W każdym kroku zliczana jest liczba pojazdów zatrzymanych na trzynastu pasach wlotowych, przy czym za pojazd należący do kolejki uznawany jest pojazd poruszający się z prędkością poniżej progu zatrzymania stosowanego przez SUMO \cite{sumo_lane_docs}. Wartości sumowane są w obrębie czterech wlotów, a następnie uśredniane po wszystkich krokach symulacji.

Metryka wyrażana jest w pojazdach, a wartości niższe są korzystniejsze. Poza wartością średnią rejestrowano również wartość maksymalną oraz wartości w podziale na poszczególne wloty, co pozwala ocenić, czy sterowanie nie poprawia sytuacji na jednym wlocie kosztem pozostałych.

\subsubsection{Przepustowość}
\label{subsubsec:metryka-przepustowosc}

Przepustowość wyznaczana jest na podstawie liczby podróży zakończonych w trakcie symulacji, czyli pojazdów, które opuściły sieć drogową. Liczba ta jest dzielona przez czas od najwcześniejszego wyjazdu do najpóźniejszego przyjazdu wśród zakończonych podróży, a następnie przeliczana na godzinę zgodnie z zależnością
\begin{equation}
    Q = \frac{N}{t_{\max} - t_{\min}} \cdot 3600,
    \label{eq:throughput}
\end{equation}
gdzie $N$ oznacza liczbę zakończonych podróży, $t_{\min}$ najwcześniejszy czas wyjazdu w tym zbiorze, a $t_{\max}$ najpóźniejszy czas przyjazdu.

Metryka wyrażana jest w pojazdach na godzinę i jest jedyną z rozpatrywanych, dla której wartości wyższe są korzystniejsze. Należy zaznaczyć, że przy zadanym z góry napływie pojazdów metryka ta służy przede wszystkim do wykrycia sytuacji, w której sterowanie nie jest w stanie obsłużyć zgłaszanego popytu.

\subsubsection{Częstotliwość przełączeń}
\label{subsubsec:metryka-przelaczenia}

Częstotliwość przełączeń opisuje, jak często zmienia się układ sygnałów na skrzyżowaniu. Za przełączenie uznawane jest przejście między dwiema różnymi konfiguracjami sygnału zielonego. Fazy żółte, stanowiące jedynie przejście między fazami zielonymi, nie są liczone jako osobne przełączenia, dzięki czemu jedna zmiana fazy nie jest zliczana dwukrotnie. Liczba przełączeń dzielona jest przez czas między pierwszym i ostatnim rekordem szeregu czasowego, a następnie przeliczana na godzinę.

Metryka wyrażana jest w przełączeniach na godzinę. W jej przypadku nie można wskazać jednoznacznie pożądanego kierunku zmian. Zbyt rzadkie przełączanie oznacza wydłużone oczekiwanie na wlotach nieobsługiwanych, natomiast zbyt częste prowadzi do narastania czasów międzyzielonych, w których żaden strumień nie korzysta ze skrzyżowania, a także obniża komfort i przewidywalność jazdy. Metryka ta pełni zatem rolę kontrolną i pozwala ocenić, czy poprawa pozostałych wskaźników nie została osiągnięta kosztem niestabilnego sterowania.

\subsection{Podział danych eksperymentalnych}
\label{subsec:podzial-danych}

Scenariusze ruchu wygenerowano przy użyciu skryptu \texttt{randomTrips.py}, dostarczanego wraz z SUMO \cite{sumo_randomtrips_docs}. Natężenie ruchu regulowano parametrem określającym średni odstęp między kolejnymi wyjazdami pojazdów, a zróżnicowanie układu tras zapewniono przez zmianę ziarna generatora liczb losowych. Wygenerowane podróże przekształcano następnie na poprawne trasy w analizowanej sieci.

Dane podzielono na trzy rozłączne zbiory: treningowy, walidacyjny i testowy. Podział ten przedstawiono na rysunku~\ref{fig:podzial-danych}.

\begin{figure}[!htbp] \centering
    \includegraphics[width=\linewidth]{podzial_danych.png}
    \caption{Podział danych eksperymentalnych oraz przeznaczenie poszczególnych zbiorów.}
    \label{fig:podzial-danych}
\end{figure}

\subsubsection{Zbiór treningowy}
\label{subsubsec:zbior-treningowy}

Przygotowano dwa warianty zbioru treningowego, różniące się sposobem rozmieszczenia czasów wyjazdu pojazdów. Celem było sprawdzenie, czy większa nieregularność napływu pojazdów w danych treningowych poprawia zdolność modelu do generalizacji.

Pierwszy wariant obejmuje scenariusze o regularnych odstępach między wyjazdami, generowane dla zadanego okresu napływu. Drugi wariant obejmuje scenariusze o zbliżonych poziomach natężenia, lecz z nieregularnymi czasami wyjazdu, uzyskanymi przez zastosowanie opcji \texttt{--random-depart} oraz dodatkowego różnicowania napływu pojazdów.

Każdy z wariantów zawiera 30 plików tras, po 10 plików dla każdego z trzech poziomów natężenia ruchu: lekkiego, średniego oraz wysokiego. Każdy scenariusz treningowy obejmuje 1800~s napływu pojazdów, natomiast epizod trwa 2100~s, aby umożliwić opróżnienie sieci. Poszczególne pliki różnią się ziarnem generatora liczb losowych. Podczas treningu kolejne epizody wykorzystują kolejne pliki tras, dzięki czemu agent styka się z różnymi układami tras oraz różnymi poziomami obciążenia.

\subsubsection{Zbiór walidacyjny}
\label{subsubsec:zbior-walidacyjny}

Zbiór walidacyjny obejmuje cztery scenariusze: trzy odpowiadające stałym poziomom natężenia ruchu, to jest lekkiemu, średniemu i wysokiemu, oraz jeden scenariusz dynamiczny, w którym poziom obciążenia zmienia się w trakcie symulacji. Każdy scenariusz walidacyjny obejmuje 3600~s napływu pojazdów.

Zbiór ten wykorzystano do podejmowania wszystkich decyzji dotyczących konfiguracji modelu, a w szczególności do:
\begin{itemize}
    \item wyboru odpowiedniej liczby kroków treningowych na podstawie oceny kolejnych checkpointów;
    \item porównania obu wariantów zbioru treningowego;
    \item doboru hiperparametrów algorytmu;
    \item oceny stabilności procesu uczenia;
    \item wyboru modelu końcowego.
\end{itemize}

\subsubsection{Zbiór testowy}
\label{subsubsec:zbior-testowy}

Zbiór testowy obejmuje cztery grupy scenariuszy, odpowiadające trzem stałym poziomom natężenia ruchu oraz scenariuszowi o zmiennym natężeniu. W każdej grupie przygotowano pięć wariantów różniących się ziarnem generatora liczb losowych, co daje łącznie 20 plików tras. Każdy scenariusz testowy obejmuje 3600~s napływu pojazdów.

Zastosowanie pięciu wariantów w obrębie każdej grupy pozwala ograniczyć wpływ pojedynczego, korzystnego lub niekorzystnego układu tras na wnioski. Wyniki przedstawiane są zatem jako wartości uśrednione wraz z miarą rozrzutu, a nie jako pojedyncze przebiegi. Te same pliki tras wykorzystano do oceny wszystkich trzech porównywanych metod sterowania.

\subsubsection{Zapobieganie przenikaniu danych}
\label{subsubsec:przenikanie-danych}

Rozdzielenie zbiorów jest istotnym elementem metodyki, ponieważ wielokrotne korzystanie z tych samych danych na etapie strojenia i oceny prowadziłoby do zawyżenia wyników i utraty wiarygodności wniosków. W eksperymencie przyjęto następujące zasady:
\begin{itemize}
    \item wszystkie decyzje dotyczące liczby kroków treningowych, wyboru zbioru treningowego oraz wartości hiperparametrów podejmowano wyłącznie na podstawie wyników na zbiorze walidacyjnym;
    \item parametry czasowe metod referencyjnych, to jest czasy trwania faz programu stałoczasowego oraz ograniczenia czasowe sterowania akomodacyjnego, wyznaczono wyłącznie na podstawie scenariuszy ze zbioru treningowego;
    \item konfigurację modelu końcowego ustalono i zamrożono przed jego pierwszą ewaluacją na zbiorze testowym;
    \item zbioru testowego nie wykorzystywano podczas treningu ani strojenia, a jego scenariusze zostały wygenerowane z użyciem ziaren niezależnych od pozostałych zbiorów;
    \item po obejrzeniu wyników testowych nie wprowadzano zmian w modelu ani w jego konfiguracji.
\end{itemize}
